% !PRINTED 338
\weberpage{338}
% running head: 338 | Elfter Abschnitt. | §. 91.

In order to compute $\varDelta_{\lambda,\mu}$ for the canonical form, one has,
in \S.\,90, (3), to replace
\[
h, \qquad m
\]
by
\[
h(\lambda-6m^2\mu), \qquad m\lambda+(h^3+2m^3)\mu
\]
and one thus finds:
\[
\varDelta_{\lambda,\mu} = -\,h(\lambda-6m^2\mu)\,[m\lambda+(h^3+2m^3)\mu]^2\,\Sigma y_1^3
+ \{h^3(\lambda-6m^2\mu)^3+2[m\lambda+(h^3+2m^3)\mu]^3\}\,y_1y_2y_3,
\]
from which, by \S.\,90, (7), (8), with regard to (2), it follows after some
simple calculations that:
\weq{5}{\begin{gathered}
L = -2\lambda^2\mu(h^3-m^3)m - \lambda\mu^2(h^6-20h^3m^3-8m^6) \\
\qquad{}+ 8\mu^3m^2(h^3-m^3)^2, \\
M = \lambda^3+12\lambda\mu^2m(h^3-m^3)+2\mu^3(h^6-20h^3m^3-8m^6).
\end{gathered}}

In this way, then, one obtains only two invariants, of degree
$4$ and degree $6$, namely, in the canonical form:
\weq{6}{m(h^3-m^3), \quad h^6-20h^3m^3-8m^6.}

The weight of these invariants is [by Bd. I, \S.\,60, (3)]
likewise 4 and 6, and if we therefore denote them in the general
form by $S$ and $T$ and set $h=1$, we find
\weq{7}{\begin{gathered}
m(1-m^3) = r^4S, \\
1-20m^3-8m^6 = r^6T.
\end{gathered}}

For $L$ and $M$ one obtains from (5) the representation:
\weq{8}{\begin{gathered}
L = -2S\lambda^2\mu - T\lambda\mu^2 + 8S^2\mu^3 \\
M = \lambda^3+12S\lambda\mu^2+2T\mu^3.
\end{gathered}}

The functions $S$ and $T$ of degree $4$ and degree $6$ were
first computed by \gesperrt{Aronhold}. We shall not set down the long
expressions here, but from now on regard these two
quantities as known and given.

From (7) a biquadratic equation can be derived
for the ratio $m^2:r^2$. One finds from it, namely:
\weq{9}{\begin{gathered}
m^5-m^8 = r^4m^4S \\
m^2-2m^5+m^8 = r^8S^2 \\
m^2-20m^5-8m^8 = r^6m^2T.
\end{gathered}}

If one eliminates $m^2$ and $m^5$ from these, one finds
\weq{10}{27m^8+18Sm^4r^4+Tm^2r^6-S^2r^8=0.}
